\scanpage{32}
1\deleted{\unsure{5}}6) \uline{Classes fondamentales} \uline{globales}

On fixe un entier $n>0$, et on se \deleted{Soit $f:X\unsure{\mathrel{-}|\longrightarrow}Y$ morphisme} borne à des faisceaux de $n$-torsion (i.e. Modules sur $A=\mathbb Z/n\mathbb Z$), et des préschémas à \unsure{car.} résiduelles premières à $n$ i.e tels que $n.1_{\mathcal O_X}$ est inversible dans $\mathcal O_X$.

Si $f:X\mathrel{-}Y$ est un morphisme lisse, on dira qu'il est \inserted{purement} de dim relative $d$ si ses fibres sont \inserted{purement} de dim $d$. On \inserted{désigne} par $T^{(n)}_{X/Y}$ ou simplement $T_{X/Y}$ s'il n'y a pas de confusion \unsure{possible|possibles}, le faisceau $\unsure{\mu_{nX}|\mu_{n_X}}^{\otimes d}$. [Si $f$ est lisse de dim \inserted{relative \unsure{év\textsuperscript{t}.|év!} variable} quelconque, alors on définit de façon évidente $T_{X/Y}$ en le posant sur les divers morceaux].

Nous \unsure{regardons|regarderons} par la suite des morphismes \uline{lisses} \uline{compactifiables} \uline{purs} \uline{de dim} $d$ \qquad $f:X\longrightarrow Y$.

Pour un tel morphisme, on peut \unsure{définir|définit} un \uline{\deleted{\unsure{prem|prem.}}\inserted{hom.} canonique}

% archive sheet number: 31
